\documentclass[11pt]{article}
\usepackage[a4paper,margin=1in]{geometry}
\usepackage{amsmath,amssymb,amsthm,mathtools,booktabs}
\usepackage{hyperref}

\title{Finite Confinement of Infinity: Residues and Equilibrium Constants for Compactifying Maps}
\author{Author Name}
\date{\today}

\theoremstyle{definition}
\newtheorem{definition}{Definition}[section]
\newtheorem{remark}[definition]{Remark}
\newtheorem{conjecture}[definition]{Conjecture}

\theoremstyle{plain}
\newtheorem{lemma}[definition]{Lemma}
\newtheorem{proposition}[definition]{Proposition}
\newtheorem{theorem}[definition]{Theorem}
\newtheorem{corollary}[definition]{Corollary}

\begin{document}
\maketitle

\begin{abstract}
We formalize two compactifying maps, 
\(f_1(x)=\frac{x}{1+x}\) and \(f_2(x)=\frac{2}{\pi}\arctan x\), and their arithmetic mean 
\(f_{\mathrm{avg}}\).
Using a fixed residue-at-infinity convention and explicit equilibrium equations, we derive finite constants from asymptotic expansions and from value--slope balance points.
We separate theorem-level results from deferred constants that still require independent defining equations.
\end{abstract}

\section{Introduction}
This article isolates a finite-confinement layer used in the \(\Omega_p\) program.
Its scope is restricted to elementary analytic statements about selected compactifying maps.
No claim about the Riemann Hypothesis or other Millennium problems is proved here.

\section{State of the Art}
Compactifications such as rational saturation and arctangent normalization are standard in analysis and numerical modeling.
The novel contribution here is only organizational: a shared residue convention and an explicit status table (proved vs.
conjectural) for constants used downstream.

\section{Notation and Definitions}
\begin{definition}[Maps on \((0,\infty)\)]
\label{def:maps}
For \(x>0\), define
\[
f_1(x)=\frac{x}{1+x}, \qquad
f_2(x)=\frac{2}{\pi}\arctan x, \qquad
f_{\mathrm{avg}}(x)=\frac{f_1(x)+f_2(x)}{2}.
\]
\end{definition}

\begin{definition}[Residue-at-infinity convention]
\label{def:res-inf}
If a meromorphic function has Laurent expansion at infinity
\[
f(z)=a_0+\frac{a_{-1}}{z}+\frac{a_{-2}}{z^2}+\cdots,
\]
we use the classical convention
\[
\operatorname{Res}_{\infty} f := -a_{-1}.
\]
\end{definition}

\begin{definition}[Compressed coefficient]
\label{def:compressed}
Define the compressed coefficient by
\[
\mathcal C_{\infty}(f):=\operatorname{Res}_{\infty} f.
\]
Hence \(\mathcal C_{\infty}(f)\) is the positive coefficient that appears in expansions of the form
\(f(x)=1-\mathcal C_{\infty}(f)x^{-1}+o(x^{-1})\).
\end{definition}

\begin{definition}[Primary equilibrium constants]
\label{def:primary}
Define \(C_1,C_3,C_5\) as the unique positive solutions of
\[
f_1(x)=f_1'(x),\qquad f_2(x)=f_2'(x),\qquad f_{\mathrm{avg}}(x)=f_{\mathrm{avg}}'(x).
\]
\end{definition}

\begin{definition}[Catalog constants with deferred equations]
\label{def:deferred}
Symbols \(C_2,C_4,C_6\) denote secondary equilibrium constants from the legacy catalog.
In this article they are recorded as placeholders; their defining equations are deferred and therefore they are not theorem-level objects here.
\end{definition}

\section{Preliminary Lemmas}
\begin{lemma}
\label{lem:f1-expansion}
For \(|z|>1\),
\[
\frac{z}{1+z}=1-\frac{1}{z}+\frac{1}{z^2}-\frac{1}{z^3}+\cdots.
\]
Hence \(\operatorname{Res}_{\infty} f_1=1\).
\end{lemma}
\begin{proof}
Use \(\frac{1}{1+z^{-1}}=1-z^{-1}+z^{-2}-\cdots\), then multiply by \(1\).
The coefficient of \(z^{-1}\) is \(-1\), so by Definition~\ref{def:res-inf},
\(\operatorname{Res}_{\infty} f_1=-(-1)=1\).
\end{proof}

\begin{lemma}
\label{lem:f2-expansion}
As \(|z|\to\infty\) in sectors avoiding branch crossing,
\[
\frac{2}{\pi}\arctan z = 1-\frac{2}{\pi}\frac{1}{z}+O\!\left(\frac{1}{z^3}\right).
\]
Hence \(\operatorname{Res}_{\infty} f_2=\frac{2}{\pi}\).
\end{lemma}
\begin{proof}
Use the asymptotic expansion \(\arctan z=\frac{\pi}{2}-z^{-1}+O(z^{-3})\).
Multiplying by \(\frac{2}{\pi}\) yields the result.
\end{proof}

\begin{lemma}
\label{lem:avg-expansion}
\[
f_{\mathrm{avg}}(z)=1-\left(\frac12+\frac1\pi\right)\frac{1}{z}+O\!\left(\frac{1}{z^2}\right),
\]
so
\[
\operatorname{Res}_{\infty}f_{\mathrm{avg}}=\frac12+\frac1\pi.
\]
\end{lemma}
\begin{proof}
Average the expansions from Lemmas~\ref{lem:f1-expansion} and \ref{lem:f2-expansion}.
\end{proof}

\begin{lemma}
\label{lem:monotone-balance}
For each \(f\in\{f_1,f_2,f_{\mathrm{avg}}\}\), the balance map
\(h_f(x):=f(x)-f'(x)\) is strictly increasing on \((0,\infty)\).
\end{lemma}
\begin{proof}
Direct differentiation gives
\(h_f'(x)=f'(x)-f''(x)\).
For \(f_1\): \(f_1'(x)=\frac{1}{(1+x)^2}\), \(f_1''(x)=-\frac{2}{(1+x)^3}\), so
\(h_{f_1}'(x)=\frac{x+3}{(1+x)^3}>0\).
For \(f_2\): \(f_2'(x)=\frac{2}{\pi(1+x^2)}\), \(f_2''(x)=-\frac{4x}{\pi(1+x^2)^2}\), so
\(h_{f_2}'(x)=\frac{2(1+x)^2}{\pi(1+x^2)^2}>0\).
For \(f_{\mathrm{avg}}\), linearity gives
\(h_{f_{\mathrm{avg}}}'=\frac{1}{2}(h_{f_1}'+h_{f_2}')>0\).
\end{proof}

\section{Main Theorems}
\begin{theorem}[Residues and compressed constants]
\label{thm:residues}
With Definitions~\ref{def:res-inf}--\ref{def:compressed},
\[
\mathcal C_{\infty}(f_1)=1,
\qquad
\mathcal C_{\infty}(f_2)=\frac{2}{\pi},
\qquad
\mathcal C_{\infty}(f_{\mathrm{avg}})=\frac12+\frac1\pi=:C_{\mathrm{avg}}.
\]
Numerically, \(C_{\mathrm{avg}}\approx 0.8183098861837907\).
\end{theorem}
\begin{proof}
Immediate from Lemmas~\ref{lem:f1-expansion}--\ref{lem:avg-expansion}.
\end{proof}

\begin{theorem}[Existence and uniqueness of primary equilibrium constants]
\label{thm:primary}
Each equation in Definition~\ref{def:primary} has exactly one positive solution.
Hence constants \(C_1,C_3,C_5\) are well-defined and unique.
\end{theorem}
\begin{proof}
By Lemma~\ref{lem:monotone-balance}, each \(h_f\) is strictly increasing.
Also, for each map, \(h_f(0^+)=-f'(0)<0\) and \(\lim_{x\to\infty}h_f(x)=1>0\).
By the intermediate value theorem, each \(h_f\) has at least one zero.
Strict monotonicity implies uniqueness.
\end{proof}

\begin{proposition}[Numerical approximations (certified by bisection)]
\label{prop:numerics}
Using deterministic bisection on sign-changing intervals:
\[
C_1\approx 0.6180339887498949,\qquad
C_3\approx 0.7472119551615677,\qquad
C_5\approx 0.6829108387479741.
\]
\end{proposition}
\begin{proof}
See computational protocol in the appendix note (separate methods file).
\end{proof}

\section{Proved vs. Conjectural}
\subsection*{Proved in this article}
\begin{itemize}
\item Residue-at-infinity values for \(f_1,f_2,f_{\mathrm{avg}}\) under a fixed convention.
\item Closed form \(C_{\mathrm{avg}}=\frac12+\frac1\pi\).
\item Existence and uniqueness of \(C_1,C_3,C_5\) (primary equilibrium constants).
\item Reproducible numerical approximations for \(C_1,C_3,C_5\).
\end{itemize}

\subsection*{Deferred / conjectural in this article}
\begin{itemize}
\item Defining equations for \(C_2,C_4,C_6\) are not fixed here.
\item Any universal-residue claim across unrelated algebraic systems.
\item Any implication for RH, \(\zeta\)-zeros, or Millennium problems.
\end{itemize}

\begin{conjecture}[Secondary equilibrium completion]
There exists a mathematically natural and non-circular definition of
\(C_2,C_4,C_6\) compatible with the primary constants and with explicit uniqueness proofs.
\end{conjecture}

\section{Limitations and Falsifiability}
This paper uses only three compactifying maps and does not claim universality.
The numerical constants are algorithm-dependent approximations, while the existence/uniqueness results are analytic.
Any future attempt to use this layer in larger claims must provide independent, non-circular bridges.

\section{Conclusion}
The article provides a rigorous finite-confinement baseline:
residue convention fixed, primary constants defined and proved unique,
and deferred constants explicitly marked as unfinished.
This supports editorial clarity for subsequent articles while preventing over-claiming.

\bibliographystyle{plain}
\begin{thebibliography}{9}
\bibitem{ahlfors} L. V. Ahlfors, \emph{Complex Analysis}.
\bibitem{rudin} W. Rudin, \emph{Real and Complex Analysis}.
\bibitem{stoer} J. Stoer and R. Bulirsch, \emph{Introduction to Numerical Analysis}.
\end{thebibliography}

\end{document}
